\documentclass[11pt,a4paper]{article}
% Estilo común. Editar aquí para mantener coherentes los tres documentos.
% Adaptación de las plantillas del Seminario Agentes Inteligentes y LLM (UNLu).
\usepackage[utf8]{inputenc}
\usepackage[T1]{fontenc}
\usepackage[spanish,es-nodecimaldot]{babel}
\usepackage[default]{sourcesans}
\usepackage{microtype}
\usepackage[a4paper,left=23mm,right=23mm,top=23mm,bottom=23mm,headheight=15pt,headsep=8mm,footskip=11mm]{geometry}
\usepackage{xcolor}
\usepackage{graphicx}
\usepackage{enumitem}
\usepackage{titlesec}
\usepackage{fancyhdr}
\usepackage{needspace}
\usepackage[skins,breakable]{tcolorbox}
\usepackage{hyperref}
\usepackage{xurl}

\definecolor{azul}{HTML}{1A465E}
\definecolor{petroleo}{HTML}{2F859A}
\definecolor{texto}{HTML}{294D60}
\definecolor{suave}{HTML}{EFF7F9}
\definecolor{borde}{HTML}{CFE2E8}
\definecolor{gris}{HTML}{586D79}

\hypersetup{colorlinks=true,urlcolor=petroleo,linkcolor=azul,citecolor=azul,
  pdfauthor={Aprendizaje Automático - Universidad Nacional de Luján}}
\urlstyle{same}
\color{texto}
\setlength{\parindent}{0pt}
\setlength{\parskip}{6pt}
\linespread{1.10}
\setlength{\emergencystretch}{1.5em}
\clubpenalty=10000
\widowpenalty=10000

% Más aire antes del título que después: el título acompaña su contenido.
\titleformat{\section}{\Large\bfseries\color{azul}}{}{0pt}{}
\titlespacing*{\section}{0pt}{24pt}{0pt}
\titleformat{\subsection}{\large\bfseries\color{azul}}{}{0pt}{}
\titlespacing*{\subsection}{0pt}{18pt}{0pt}
\setlist[itemize]{leftmargin=16pt,itemsep=4pt,topsep=4pt,parsep=0pt,label=\textcolor{petroleo}{\textbullet}}
\setlist[enumerate]{leftmargin=20pt,itemsep=6pt,topsep=4pt,parsep=0pt,label=\textcolor{petroleo}{\bfseries\arabic*.}}

\pagestyle{fancy}
\fancyhf{}
\fancyhead[L]{\small\color{gris}Universidad Nacional de Luján}
\fancyhead[R]{\small\color{gris}Aprendizaje Automático · 2026}
\fancyfoot[L]{\small\color{gris}\pieDocumento}
\fancyfoot[R]{\small\color{petroleo}\thepage}
\renewcommand{\headrulewidth}{0.35pt}
\renewcommand{\footrulewidth}{0pt}
\newcommand{\pieDocumento}{Trabajo práctico integrador}

\newcommand{\cabecera}[3]{%
  {\small\bfseries\color{petroleo}\MakeUppercase{#1}\par}
  \vspace{4pt}
  {\fontsize{24}{28}\selectfont\bfseries\color{azul}#2\par}
  \vspace{7pt}
  {\large\color{gris}#3\par}
  \vspace{6pt}}

\newtcolorbox{destacado}[1]{enhanced,breakable,colback=suave,colframe=borde,
  boxrule=0.5pt,arc=2mm,left=12pt,right=12pt,top=9pt,bottom=9pt,
  borderline west={2pt}{0pt}{petroleo},before skip=12pt,after skip=10pt,
  title={#1},colbacktitle=suave,coltitle=azul,fonttitle=\bfseries,
  attach title to upper={\par\smallskip},fontupper=\normalsize}

\newcommand{\orientacion}[1]{{\color{gris}#1}\par}
\newcommand{\campo}[1]{\textcolor{petroleo}{\textit{[#1]}}}
% Encabezado compacto para los ejemplos consecutivos de la consigna.
\newcommand{\lineadesarrollo}[1]{\par\vspace{12pt}\noindent{\large\bfseries\color{azul}#1\par}\nopagebreak\vspace{-2pt}}

\renewcommand{\pieDocumento}{Plantilla de informe · Trabajo práctico integrador}
\hypersetup{pdftitle={Plantilla de informe - Aprendizaje Automático}}

% Adaptación del short paper del Seminario Agentes Inteligentes y LLM.
% Se conserva la secuencia problema -> metodología -> implementación -> evidencia.
% Se simplifican el estado del arte y la arquitectura y se centra el informe
% en los datos reales, la evaluación y el uso de la solución.
% El informe acompaña obligatoriamente al proyecto; usar LaTeX es opcional.

\begin{document}
\cabecera{Trabajo práctico integrador}{\campo{Título del proyecto}}{Aprendizaje Automático · Universidad Nacional de Luján · 2026}
\textbf{Estudiante/s:} \campo{Nombre y apellido}\par
\textbf{Repositorio:} \campo{Enlace a GitHub o plataforma equivalente}\par

\begin{destacado}{Cómo utilizar esta plantilla}
Reemplacen los campos y las orientaciones por el contenido del proyecto. La estructura es una guía adaptable: ajusten los apartados al enfoque y la extensión al trabajo realizado. El informe en PDF acompaña obligatoriamente al proyecto y debe estar enlazado desde el README; no reemplaza la solución reproducible. Usar LaTeX es opcional.
\end{destacado}

\section*{Resumen}
\orientacion{Presenten el problema real, los datos utilizados, el enfoque de IA y la principal evidencia obtenida. Escriban el resumen al finalizar el trabajo y distingan lo implementado de aquello que proponen como continuación.}
\campo{Resumen del trabajo.}

\section{Situación, objetivo y alcance}
\orientacion{Describan dónde aparece el problema, a quién afecta y qué tarea se busca apoyar. Formulen un objetivo concreto y expliquen cómo reconocerían un resultado útil. Delimiten qué parte de la situación resuelve el prototipo y qué queda fuera de su alcance.}
\campo{Contexto real, destinatario, objetivo y alcance.}

\section{Datos y representación del problema}
\orientacion{Identifiquen las fuentes reales, la fecha o versión del conjunto utilizado y las condiciones de acceso. Describan los atributos, clases, imágenes, textos o registros relevantes, los problemas de calidad encontrados y la preparación realizada. Si trabajan con reglas, estados o acciones, expliquen su significado y la evidencia del dominio que sustenta esa representación.}
\campo{Fuentes, características, preparación y representación.}

\newpage
\section{Enfoque y metodología}
\orientacion{Justifiquen el enfoque y su relación con los conceptos del curso. Si requiere aprendizaje, describan cómo entrenaron o ajustaron el modelo: datos, arquitectura o técnica, función de pérdida, hiperparámetros relevantes y criterio para elegir la configuración final. Analicen la evolución del entrenamiento cuando corresponda. Si utilizaron un modelo preentrenado sin ajustarlo, justifiquen esa decisión. En soluciones simbólicas o de búsqueda, expliquen las reglas, la representación y los criterios de decisión.}
\campo{Técnica elegida, entrenamiento o ajuste cuando corresponda, justificación y principales decisiones.}

\subsection{Diseño de la evaluación}
\orientacion{Indiquen qué van a medir u observar y por qué permite evaluar el objetivo. Si existe entrenamiento, expliquen la separación de los datos y cómo evitan usar información de prueba durante el ajuste. En otras soluciones, describan los casos de prueba y sus resultados esperados. Si comparan con una referencia, expliciten cuál es y por qué resulta pertinente.}
\campo{Criterios, datos de evaluación, casos de prueba y referencia de comparación.}

\section{Implementación y funcionamiento}
\orientacion{Expliquen qué construyeron y cómo se utiliza. Presenten las entradas, las transformaciones principales y las salidas. Un esquema sencillo puede ayudar, pero debe estar acompañado por una explicación. No es necesario copiar el código completo: indiquen dónde está en el repositorio y cuáles son las decisiones que ayudan a comprenderlo.}
\campo{Flujo de funcionamiento y decisiones de implementación.}

\subsection{Reproducción}
\orientacion{Indiquen el entorno, las dependencias, cómo obtener los datos y el orden de ejecución. Pueden remitir al README para el detalle operativo. Si utilizan modelos o servicios externos, identifíquenlos e indiquen qué se necesita para reproducir los pasos.}
\campo{Instrucciones o referencia a la sección correspondiente del README.}

\section{Resultados y análisis de errores}
\orientacion{Presenten evidencia de la versión ejecutada: métricas, tablas, visualizaciones o ejemplos analizados. Expliquen qué muestran y cómo se relacionan con el objetivo. Incluyan aciertos y fallos, condiciones en las que la solución funciona peor y límites de las conclusiones. No completen este apartado con resultados hipotéticos.}
\campo{Resultados reales del prototipo e interpretación.}

\newpage
\section{Uso de la solución en el contexto real}
\orientacion{Expliquen quién utilizaría el modelo o sistema, qué entrada aportaría, qué recibiría y qué haría con la salida. Distingan la decisión del modelo de la decisión de la persona. Señalen qué revisión humana sería necesaria y qué faltaría para integrar el prototipo a la situación real. No es obligatorio desplegar una aplicación para desarrollar esta reflexión.}
\campo{Escenario de uso, intervención prevista y condiciones de implementación.}

\section{Limitaciones, conclusiones y mejoras}
\orientacion{Retomen el objetivo y expliquen qué lograron demostrar. Consideren límites de los datos y de la evaluación, sesgos, errores, costos o condiciones de uso relevantes. Propongan mejoras vinculadas con la evidencia obtenida; eviten presentar como resuelto aquello que no evaluaron.}
\campo{Conclusiones sustentadas, limitaciones y mejoras.}

\section*{Declaración de uso de IA}
\orientacion{Un disclaimer o declaración de uso de IA es una nota de transparencia sobre la asistencia recibida durante la elaboración del trabajo. Identifiquen las herramientas o modelos utilizados (y la versión, si la conocen), las tareas en las que intervinieron, los aportes incorporados y cómo los revisaron o verificaron. No es una exención de responsabilidad: deben comprender y poder explicar lo entregado. Los modelos que integran la solución se describen además en la metodología. Si no utilizaron IA como asistencia, indíquenlo expresamente.}
\campo{Declaración que refleje el uso real de herramientas de IA.}

\begin{destacado}{Ejemplo para adaptar, no para copiar automáticamente}
Utilizamos [herramienta/modelo] para consultar errores de código y revisar la redacción. Incorporamos sugerencias en [partes del trabajo], luego de comprobar su funcionamiento mediante [pruebas] y contrastar las afirmaciones con [fuentes]. Las decisiones metodológicas y el análisis de los resultados fueron realizados por el equipo.
\end{destacado}

\section*{Fuentes y recursos utilizados}
\orientacion{Identifiquen los datos, libros, artículos, modelos, bibliotecas y fragmentos de código utilizados. Incluyan referencias o enlaces que permitan localizar los recursos. Las referencias deben corresponder a materiales efectivamente empleados en el trabajo.}
\begin{itemize}
\item \campo{Fuente de datos: responsable, nombre, versión o fecha y enlace.}
\item \campo{Fuente conceptual o técnica utilizada.}
\item \campo{Modelo, biblioteca o código reutilizado, cuando corresponda.}
\end{itemize}

% Si prefieren citas bibliográficas formales, pueden reemplazar el apartado
% anterior por thebibliography o por una bibliografía BibTeX propia.
% Ejemplo para insertar una figura real (requiere el archivo en el proyecto):
% \begin{figure}[htbp]
%   \centering
%   \includegraphics[width=0.85\linewidth]{figuras/resultado.png}
%   \caption{Descripción del resultado observado.}
%   \label{fig:resultado}
% \end{figure}

\end{document}
